\section{Percabangan: if, else if, else, dan switch}

Percabangan memungkinkan program mengeksekusi kode yang berbeda berdasarkan kondisi. Statement \code{if} mengevaluasi ekspresi boolean; jika \code{true}, blok kode di dalamnya dieksekusi. \code{else if} menambahkan kondisi tambahan, dan \code{else} menangani kasus selain semua kondisi di atasnya. JavaScript juga menyediakan \code{switch} untuk membandingkan satu nilai dengan banyak kemungkinan \cite{jsInfo}.

\begin{webcode}[language=JavaScript, caption={if/else if/else dan switch}]
// if/else if/else
const score = 85;
let grade;
if (score >= 90) {
  grade = "A";
} else if (score >= 80) {
  grade = "B";
} else if (score >= 70) {
  grade = "C";
} else {
  grade = "D";
}
console.log(grade); // "B"

// switch (gunakan === untuk perbandingan)
const day = "Senin";
switch (day) {
  case "Senin": case "Selasa":
    console.log("Kerja"); break;
  case "Sabtu": case "Minggu":
    console.log("Libur"); break;
  default:
    console.log("Kerja juga");
}
\end{webcode}

Operator ternary (\code{kondisi ? nilaiTrue : nilaiFalse}) adalah shorthand untuk \code{if/else} sederhana: \code{const status = age >= 17 ? "dewasa" : "anak";}. Cocok untuk satu ekspresi; hindari nesting ternary yang membuat kode sulit dibaca \cite{mdnJsGuide}.

\begin{konsep}
\textbf{switch vs if/else.} Gunakan \code{switch} untuk membandingkan satu variabel dengan banyak nilai konstan (kasus diskret). Gunakan \code{if/else} untuk kondisi rentang, perbandingan kompleks, atau kondisi yang tidak terkait dengan satu variabel. Jangan lupa \code{break} setelah setiap \code{case}---tanpanya, eksekusi ``jatuh'' ke case berikutnya (fall-through), yang kadang disengaja tetapi sering menjadi bug \cite{jsInfo}.
\end{konsep}

